\subsection*{Correctness of the Greedy Algorithm for the Interval Scheduling Problem}

We now show that, the greedy algorithm of picking the interval with earliest finish-time 
followed by removing the intervals conflicting with it is optimal.
Let $X$ be the set of intervals picked by this greedy algorithm.
First, all intervals in $X$ do not conflict with each other. %compatible~(even if using the first 3 strategies),
This is obvious, as right after picking an interval all conflicting ones will be immediately removed;
so, any newly added interval will not conflict with ones already picked.

Now we show that the algorithm is optimal, i.e., $|X|$ is maximized.
Assume that $X = \{x_1, x_2, \cdots, x_k\}\subseteq\{1, 2, \cdots, n\}$.
We also assume that these $k$ intervals are sorted by their finish-time, i.e., $f(x_1) \le f(x_2) \le \cdots \le f(x_k)$.
Assume that $Y = \{y_1, y_2, \cdots, y_m\}$ be an optimal solution~(with $m$ intervals).
We also assume that these $m$ intervals are sorted by their finish-time, i.e., $f(y_1) \le f(y_2) \le \cdots \le f(y_m)$.

Clearly $k \le m$, as we assume $Y$ is the optimal solution.
The ultimate goal is to prove that $k = m$, i.e., the greedy algorithm always finds the
same number of intervals with the optimal solution.
In order to achieve it, we first prove a more informative fact,
which shows that the greedy solution always \emph{stays ahead} of the optimal solution. Specifically,
the finish-time of the $i$-th interval in the greedy solution
is always no latter than the finish-time of the $i$-th interval of the optimal solution,
formally stated below.

\begin{fact} \label{int-fact1}
For any $1 \le i \le k$, we have $f(x_i) \le f(x_i^*)$.
\end{fact}

\emph{Proof.} We prove it by induction. The base case
is to show $f(x_1) \le f(y_1)$. This is obvious,
as the first interval picked by the greedy algorithm 
is always the one with smallest finish-time.

We now show the inductive step. Suppose $f(x_j) \le f(y_j)$ for all $1\le j \le i$.
We aim to prove $f(x_{i+1}) \le f(y_{i+1})$. 
See Figure~\ref{fig:inductive}.
We can write $f(x_i) \le f(y_i) \le s(y_{i+1})$, 
   where the first inequality comes from the inductive assumption
   and the second one comes from the fact that intervals in the optimal solution do not conflict.
This shows that, $y_{i+1}$ does not conflict with $x_i$,
and hence $y_{i+1}$ does not conflict with any of these $\{x_1, x_2, \cdots, x_i\}$.
Therefore, in the greedy algorithm, after picking $x_i$, the interval $y_{i+1}$ remains~(i.e., not removed).
In other words, when the greedy algorithm picks the $(i+1)$-th interval, 
interval $y_{i+1}$ is one of the available candidates.
As the greedy algorithm always pick the one with earliest finish-time~(among all available ones), 
either $y_{i+1}$ is picked, or another one with even smaller finish-time~(than $y_{i+1}$) is picked;
in either case, we have $f(x_{i+1}) \le f(y_{i+1})$, as desired. \qed

\begin{figure}[h]
\centering{\input{inductive}}
\caption{Illustrating of proving Fact~\ref{int-fact1}.}
\label{fig:inductive}
\end{figure}

\begin{fact} \label{fact2}
We have $k = m$.
\end{fact}

\emph{Proof.} We prove it by contradiction.
See Figure~\ref{fig:fact2}.  Suppose conversely that $m \ge k + 1$.
Using the same reasong in proving Fact~\ref{int-fact1},
we know that in the greedy algorithm after picking $\{x_1, x_2, \cdots, x_k\}$,
$y_{k+1}$ remains alive, as it does not overlap with any of these $k$ intervals that are already picked.
Hence, the greedy algorithm will not terminate and will continue to pick at least one more interval.
This contradicts to the assumption that the greedy solution includes $k$ intervals. \qed

\begin{figure}[h]
\centering{\input{fact2}}
\caption{Illustrating of proving Fact~\ref{fact2}.}
\label{fig:fact2}
\end{figure}


\subsection*{The Complete Algorithm}

One iteration of the greedy algorithm requires to remove the intervals that conflict with the currently picked one.
The naive approach is to compare with all remaining intervals,
which takes linear time, resulting in an overall $O(n^2)$ time algorithm.
%with earliest finish-time that is compatible with existing ones~(i.e., those in $X$)?

Can we do it faster?
We can sort all intervals by their finish-time, examine them one-by-one following this order, and include it when it does not
conflict with those already picked~(denoted as $X$). %---a similar approach is used in the Kruskal's algorithm.
Note that, this means we delay the removal: instead of removing the conflicting intervals immediately, we remove
it at the time of examining it, and actually remove/skip it if it conflicts with one of the added intervals.

The next question is, when examining the current interval, 
how to decide if it conflicts with any one in $X$?
Again we can do linear scan but it will take linear time.
%A naive solution is to compare it with all intervals in $X$, but this will take linear time and the resulting algorithm will take $O(n^2)$ time.
%In fact, we can do it faster. 
In fact, we only need to check if it conflicts with the \emph{last interval added to $X$}.
This is because, the current interval does not conflict with all in $X$ 
is equivalent to that its start-time is no sooner than the finish-time of all intervals in $X$,
but since all intervals are now sorted by their finish-time,
this becomes equivalent to that the start-time of the current interval is no sooner than the 
finish-time of the last interval added to $X$.

In implementation, we simply maintain the last interval that is added to $X$,
and use a single comparison to decide if the current one should be picked or not.
The pseudo-code is given below.
The for-loop takes linear time. The entire algorithm takes $O(n\log n)$ time, as it is dominated by sorting.
An example of running this algorithm follows.

\begin{algorithmblock}[0.8\textwidth]
	\aaA {9}{Algorithm greedy~(intervals $\{(s(i), f(i))\}, 1\le i \le n$)}\xxx
	\aab {sort $A$ by their finish-time;}\xxx
	\aab {init $X=\{1\}$;}\xxx
	\aab {let $j=1$;\ \ // maintain the last interval added to $X$}\xxx
	\aaB {4}{for $i = 2 \to n$}\xxx
	\aaC {3}{if~($s(i) \ge f(j)$)}\xxx
	\aad {pick $i$ and add it to $X$;}\xxx
	\aad {$j = i$;}\xxx
	\aab {return $X$;}\xxx
\end{algorithmblock}

\vspace*{-0.8cm}

\begin{figure}[h]
\centering{\input{jump}}
\caption{An example of running above algorithm.}
\label{fig:jump}
\end{figure}



